\documentclass[12pt,a4paper,UTF8,fontset=none]{ctexart}
\usepackage{course-style}

% Change TEXT only. Keep the visual design in course-style.sty.
\renewcommand{\CourseSubject}{【学科名称】}
\renewcommand{\CourseTopic}{【章节主题】}
\renewcommand{\CourseEyebrow}{从第一性原理学习・【学科名称】专题课}
\renewcommand{\CourseTitleTop}{从第一性原理学会}
\renewcommand{\CourseTitleBottom}{【课程标题】}
\renewcommand{\CourseSubtitle}{从“完全不会”到能独立解决典型题}
\renewcommand{\CourseRouteA}{先修概念\\与直觉}
\renewcommand{\CourseRouteB}{定理推导\\与反例}
\renewcommand{\CourseRouteC}{典型例题\\与迁移}
\renewcommand{\CoursePromise}{\textbf{三项承诺：}每个重要概念先解释必要性；每个方法说明成立条件；每个例题从观察到证明再到迁移。}
\renewcommand{\CourseFooter}{核心概念・第一性原理推导・完整例题・分层练习}
\renewcommand{\CourseEdition}{学习版 · 2026 年}

\begin{document}
\CourseCover

\section*{先说读法：这份课怎样学？}
这不是只给答案的资料。学完后要能回答三个问题：\textbf{它是什么？为什么这样？我什么时候自己能用？}

\begin{takeawaybox}
\textbf{学习目标：}从零理解本章核心概念、能够独立证明或计算至少两类典型题，并解释常见错误发生的原因。
\end{takeawaybox}

\tableofcontents
\clearpage

\section{地基：必须补齐的先修知识}
请从最简单的例子开始，不要默认读者已经会任何未定义符号。

\subsection{概念从哪里来？}
用一个看得见、算得出的具体问题引出概念，再给出严格定义。必须说明\textbf{定义里的每个条件为什么需要}。

\begin{whybox}
为什么这个概念不是换个名字？比较去掉条件前后的例子，展示概念究竟解决什么问题。
\end{whybox}

\section{工具：从基础推出定理，而不是背结论}
陈述适用条件，逐步推导，注明直观解释与严格证明的区别，并指出不能使用的反例。

\section{例题：从读题到写下完整证明}
\begin{questionbox}
这里放从用户照片/题目/章节选出的第一道典型例题，忠实保留已知和求证。
\end{questionbox}

\stephead{第 1 步：翻译题目。} 把题目中每个条件换成具体含义。

\stephead{第 2 步：识别信号。} 解释为什么选这个工具而不是另一个工具。

\begin{trybox}
先遮住下面的内容，自己想两分钟：你能推到哪一步？
\end{trybox}

\stephead{第 3 步：完整推导。} 每次变形说明依据，不能从结论倒推代替证明。

\begin{takeawaybox}
\textbf{可迁移的方法：}说明换一个数字、换一个条件时哪些步骤仍然成立，哪些不再成立。
\end{takeawaybox}

\section{分层练习：先独立做，再看答案}
\begin{enumerate}[label=\textbf{A\arabic*.}]
\item 【基础】：检验定义是否理解。
\item 【模仿】：复现一条推导链。
\item 【迁移】：条件略变，判断原方法是否仍适用。
\end{enumerate}

\section{参考答案与错误辨析}
给每个练习写\textbf{关键步骤、为什么、结果、常见误区}，不可只贴答案。

\section{一页回顾：我真的学会了吗？}
\begin{takeawaybox}
不用翻书，用自己的话解释本章的核心定理；给出一个反例；独立完成一道未看过的迁移题。
\end{takeawaybox}

\end{document}
